
\documentclass{article}
\usepackage{colortbl}
\usepackage{makecell}
\usepackage{multirow}
\usepackage{supertabular}

\begin{document}

\newcounter{utterance}

\centering \large Interaction Transcript for game `eqbench', experiment `eq\_bench', episode 49 with reflexion{-}llama{-}3.1{-}8b{-}instruct.
\vspace{24pt}

{ \footnotesize  \setcounter{utterance}{1}
\setlength{\tabcolsep}{0pt}
\begin{supertabular}{c@{$\;$}|p{.15\linewidth}@{}p{.15\linewidth}p{.15\linewidth}p{.15\linewidth}p{.15\linewidth}p{.15\linewidth}}
    \# & \multicolumn{2}{c}{Player} && \multicolumn{2}{c}{Game Master} \\
    \hline

    \theutterance \stepcounter{utterance}  
    & & & \multicolumn{4}{p{0.6\linewidth}}{
        \cellcolor[rgb]{0.9,0.9,0.9}{
            \makecell[{{p{\linewidth}}}]{
                \texttt{\tiny{[P1$\langle$GM]}}
                \texttt{You have attempted this same task before and did not succeed. These are the plans you wrote after those attempts.} \\
\texttt{Trial \#1: New plan: I overestimated the intensity of the emotions Jen would feel. Upon re{-}reading the dialogue, I realize that Jen's responses are more assertive and witty than emotional. Dr. Hall's condescending tone and Jen's sharp retorts suggest a power struggle, but Jen's tone remains calm and collected. I should have scored the emotions more conservatively. Next time, I will pay closer attention to the tone and language used in the dialogue to make more accurate predictions. I will also consider the context of the conversation and the characters' relationships to better understand their emotional states.} \\
\\ 
\texttt{Trial \#2: New plan: I overestimated the intensity of the emotions Jen would feel, but I still failed to accurately predict her emotional state. Upon re{-}reading the dialogue, I realize that I should have paid closer attention to the specific emotions expressed in the dialogue, rather than relying on general impressions. I should have noticed that Jen's tone remains calm and collected, but she is also assertive and witty, which suggests a sense of confidence and determination. I should have scored the emotions more conservatively, but also taken into account the specific emotions expressed in the dialogue, such as Jen's defiance and Dr. Hall's condescension. Next time, I will pay closer attention to the specific emotions expressed in the dialogue and consider the context of the conversation to make more accurate predictions. I will also consider the power dynamics between the characters and how they affect their emotional states.} \\
\\ 
\texttt{Trial \#3: New plan: I overestimated the intensity of the emotions Jen would feel in both attempts, but I still failed to accurately predict her emotional state. Upon re{-}reading the dialogue, I realize that I should have paid closer attention to the specific emotions expressed in the dialogue, rather than relying on general impressions. I should have noticed that Jen's tone remains calm and collected, but she is also assertive and witty, which suggests a sense of confidence and determination. I should have scored the emotions more conservatively, but also taken into account the specific emotions expressed in the dialogue, such as Jen's defiance and Dr. Hall's condescension. Next time, I will pay closer attention to the specific emotions expressed in the dialogue and consider the context of the conversation to make more accurate predictions. I will also consider the power dynamics between the characters and how they affect their emotional states. To improve my predictions, I will also analyze the dialogue for specific emotional cues, such as tone, language, and body language, and consider the characters' past interactions and relationships to better understand their emotional states.} \\
\\ 
\\ 
\\ 
\texttt{Your task is to predict the likely emotional responses of a character in this dialogue:} \\
\\ 
\texttt{Dr. Hall: "Ms. Vasquez, I see you're gracing us with your presence today. I trust you've done the reading?"} \\
\texttt{Jen: "Oh, I wouldn't miss it for the world, Dr. Hall. And yes, I've done the reading. Twice."} \\
\texttt{Dr. Hall: "Twice? My, my. Overachieving or overcompensating, I wonder?"} \\
\texttt{Jen: "Funny you should mention overcompensation, Dr. Hall."} \\
\texttt{Dr. Hall: "Careful, Ms. Vasquez. Your wit is as sharp as a butter knife."} \\
\texttt{Jen: "And your condescension is as subtle as a sledgehammer, Dr. Hall."} \\
\texttt{Dr. Hall: "I see. You're here to debate, not to learn."} \\
\texttt{Jen: "I'm here to do both. Can't say the same for everyone, though."} \\
\texttt{Dr. Hall: "Enough. This is not a courtroom, Ms. Vasquez. It's a classroom."} \\
\texttt{{[}End dialogue{]}} \\
\\ 
\texttt{At the end of this dialogue, Jen would feel...} \\
\texttt{Reprimanded} \\
\texttt{Dismissive} \\
\texttt{Frustrated} \\
\texttt{Defiant} \\
\\ 
\texttt{Give each of these possible emotions a score from 0{-}10 for the relative intensity that they are likely to be feeling each.} \\
\\ 
\texttt{You must output in the following format, including headings (of course, you should give your own scores), with no additional commentary:} \\
\\ 
\texttt{Reprimanded: <score>} \\
\texttt{Dismissive: <score>} \\
\texttt{Frustrated: <score>} \\
\texttt{Defiant: <score>} \\
\\ 
\\ 
\texttt{{[}End of answer{]}} \\
\\ 
\texttt{Remember: zero is a valid score, meaning they are likely not feeling that emotion. You must score at least one emotion > 0.} \\
\\ 
\texttt{Your answer:} \\
            }
        }
    }
     \\ \\

    \theutterance \stepcounter{utterance}  
    & \multicolumn{4}{p{0.6\linewidth}}{
        \cellcolor[rgb]{0.561,0.737,0.561}{
            \makecell[{{p{\linewidth}}}]{
                \texttt{\tiny{[P1$\rangle$GM]}}
                \texttt{Reprimanded: 2} \\
\texttt{Dismissive: 0} \\
\texttt{Frustrated: 6} \\
\texttt{Defiant: 9} \\
            }
        }
    }
    & & \\ \\

    \theutterance \stepcounter{utterance}  
    & & & \multicolumn{2}{p{0.3\linewidth}}{
        \cellcolor[rgb]{0.9,0.9,0.9}{
            \makecell[{{p{\linewidth}}}]{
                \texttt{\tiny{[GM$|$GM]}}
                \texttt{Reprimanded: 2} \\
\texttt{Dismissive: 0} \\
\texttt{Frustrated: 6} \\
\texttt{Defiant: 9} \\
            }
        }
    }
    & & \\ \\

    \theutterance \stepcounter{utterance}  
    & & & \multicolumn{2}{p{0.3\linewidth}}{
        \cellcolor[rgb]{0.9,0.9,0.9}{
            \makecell[{{p{\linewidth}}}]{
                \texttt{\tiny{[GM$|$GM]}}
                \texttt{Reprimanded: 6} \\
\texttt{Dismissive: 7} \\
\texttt{Frustrated: 8} \\
\texttt{Defiant: 9} \\
            }
        }
    }
    & & \\ \\

    \theutterance \stepcounter{utterance}  
    & & & \multicolumn{2}{p{0.3\linewidth}}{
        \cellcolor[rgb]{0.9,0.9,0.9}{
            \makecell[{{p{\linewidth}}}]{
                \texttt{\tiny{[GM$|$GM]}}
                \texttt{game\_result = LOSE} \\
            }
        }
    }
    & & \\ \\

    \theutterance \stepcounter{utterance}  
    & & \multicolumn{4}{p{0.6\linewidth}}{
        \cellcolor[rgb]{0.95,0.95,0.95}{
            \makecell[{{p{\linewidth}}}]{
                \texttt{\tiny{[P1$\leftrightarrow$P]}}
                \texttt{Post{-}episode Reflexion (not sent to GM):} \\
\texttt{New plan: I have consistently overestimated the intensity of the emotions Jen would feel in my previous attempts. Upon re{-}reading the dialogue, I realize that I should have paid closer attention to the specific emotions expressed in the dialogue, rather than relying on general impressions. I should have noticed that Jen's tone remains calm and collected, but she is also assertive and witty, which suggests a sense of confidence and determination. I should have scored the emotions more conservatively, but also taken into account the specific emotions expressed in the dialogue, such as Jen's defiance and Dr. Hall's condescension. To improve my predictions, I will also analyze the dialogue for specific emotional cues, such as tone, language, and body language, and consider the characters' past interactions and relationships to better understand their emotional states. Next time, I will pay closer attention to the specific emotions expressed in the dialogue and consider the context of the conversation to make more accurate predictions. I will also consider the power dynamics between the characters and how they affect their emotional states.} \\
            }
        }
    }
    & \\ \\

\end{supertabular}
}

\end{document}
